\par\SourceLeaf{144}
Eine feine Kritik des Wortlauts könnte vielleicht Gründe finden, die für die erste Hypothese sprächen. Dabei will ich mich aber umso weniger aufhalten, da ich selbst diese Hypothese nicht für zulässig halten kann, und zwar hauptsächlich deshalb, weil P. sonst mit seinen eigenen Grundsätzen (oben S. [133]) in Widerspruch stünde. Ich glaube vielmehr, dass er, verleitet durch das \textsc{Kritter}sche Gutachten — oder, wie man auch sagen kann und wie ich bald ausführlich zeigen werde, durch seine unrichtige Auffassung dieses Gutachtens — die Zahl von 15 Pensionen als eine unübersteigliche oder fast unübersteigliche Schranke betrachtet habe, unter deren Schutz er seinen Plan mit voller Sicherheit machen könne.

Einen Widerschein ähnlichen Vertrauens finde ich in E.s Votum. Es zeichnet sich dadurch aus, dass es als einziges die fatalistische Zahl 15 erwähnt; sein wesentlicher Inhalt, soweit er hierher gehört, ist folgender.

Die beiden Fragen F.s, 1 und 2, sagt E., müssten allein anhand des wirklichen Vermögens und seines Ertrags im Vergleich mit \textit{den erprobten}, von H. P. sehr einleuchtend vorgelegten Grundsätzen entschieden werden. Es könne sein, dass die verlängerte Dauer der Waisenpensionen dazu führe, dass die höchste Witwenzahl von 15 überschritten werden müsse. Das dürfe aber, weil es mit der Sicherheit der Kasse unverträglich sei, keinesfalls je geschehen. Er halte es deshalb für ratsam, die Verlängerung der Waisenpensionen nur mit der ausdrücklichen Beschränkung zu bewilligen, \textit{so lange die Anzahl der Pensionen dadurch nicht über 15 gesteigert werde}. Diese Einschränkung scheine umso notwendiger, weil sich die Berechnung der Verhältnisse der Kinder nicht auf eine \textit{so gewisse} Erfahrung stützen könne wie die beobachtete Proportion der Witwen. Ich verstehe dies so: E., dessen Auffassung von P.s Plan offenbar derjenigen von \textsc{Sch.} genau entgegengesetzt ist, halte es nach der \textsc{Kritter}-P.schen Theorie zwar für gewiss, dass die Anzahl der Witwen niemals über die berechnete Zahl (also 13) hinausgehen könne. Es sei aber nicht ebenso gewiss, dass nicht mehr als 2 Waisenpensionen hinzukommen könnten. Für den Fall, dass dies doch geschehe und dadurch die Gesamtzahl der Pensionen über 15 getrieben würde, müsse die Beschränkung der Waisenpensionen vorbehalten bleiben.

Ähnliche Bedenken, dass die Anzahl der Waisenpensionen zu niedrig angesetzt sein könnte, hatten auch einige andere Abstimmende geäußert. Ihnen ist jedoch weder im abschließenden Beschluss gefolgt worden, noch haben sie sich bisher durch die Erfahrung bestätigt. Tatsächlich haben seit der Errichtung der Witwenkasse bis jetzt noch niemals mehr als zwei Waisenpensionen gleichzeitig bestanden. Dabei will ich aber nicht unerwähnt lassen, dass während eines kurzen Teils des Jahres 1776 die Anzahl auf 3 gestiegen wäre, wenn die Erstreckung der Waisenpensionen bis zum vollendeten \(20^{\text{sten}}\) Lebensjahr schon damals gegolten hätte.

Ganz anders verhält es sich jedoch mit der \textit{Witwenzahl}. Die vermeintliche Gewissheit, dass diese nicht über 13 steigen könne, ist durch die neueren Erfahrungen zerstört worden: Sie ist bereits einmal auf 22 gestiegen, und sogar der Durchschnitt während der letzten 8 Jahre hat etwas über 19 betragen. Schon lange vor 1794 war die Anzahl der Witwen, ebenfalls ohne die Waisenpensionen mitzuzählen, einmal auf 14 gestiegen und hatte während eines fünfjährigen Zeitraums (1774—1779) den Durchschnitt von 13 gehalten. Diese Tatsache, die wohl dem Kurator, freilich aber nicht den abstimmenden Professoren bekannt sein konnte, hätte schon damals die Unrichtigkeit der Annahme von P. und E. beweisen können, wenn man sie aus dem allein zulässigen, oben S. [130] Z. [4] angedeuteten Gesichtspunkt betrachtet hätte.

Da aber die Annahme der Zahl 13 als höchste Witwenzahl auf dem \textsc{Kritter}schen Gutachten beruht, in dem die Hälfte der Zahl der bestehenden Ehen als maßgebend für das Maximum der Witwenzahl aufgestellt wird, scheint der Schluss natürlich, dass darin Fehler sein könnten. Ich bin deshalb an dem Punkt angelangt, an dem ich dieses Gutachten selbst der Kritik unterziehen muss.
\par\SourceLeaf{145}\noindent
Um diese vollkommen verständlich zu machen, muss ich einige Entwicklungen vorausschicken. Mir dabei zu folgen, könnte mancher vielleicht beschwerlich finden, wenn noch nicht von vornherein abzusehen ist, worauf sie hinauslaufen werden. Es wird deshalb, scheint mir, angemessen sein, die hauptsächlichen Einwände gegen das \textsc{Kritter}sche Gutachten und die daraus gezogenen Folgerungen gleich hier an die Spitze zu stellen.
\begin{list}{}{\setlength{\leftmargin}{2em}\setlength{\labelwidth}{1.5em}\setlength{\labelsep}{.5em}\setlength{\itemsep}{0pt}\setlength{\parsep}{0pt}}
\item[1)] Ob das Verhältnis von 1 Witwe zu 2 bestehenden Ehen auf die hiesige Professorenwitwenkasse anwendbar ist, steht nicht fest. Vielmehr ist es, wie schon oben S. [133] Z. [5] bemerkt, wahrscheinlich, dass man nach den Verhältnissen dieser Kasse mit etwas mehr Witwen rechnen müsse.
\item[2)] \textsc{Kritter}s Behauptung [Z. 10 des Abdrucks], auch wenn sie unabhängig von dem vorausgesetzten Verhältnis \(2:1\) vorgetragen wird, nämlich man könne annehmen, dass ein Durchschnitt von 10 Jahren immer schon genüge, das wahre, für die Verhältnisse der Kasse gültige Verhältnis sehr nahe anzugeben, ist durchaus falsch. Die Unrichtigkeit dieses Satzes lässt sich auch ohne Berufung auf Erfahrungen bereits aus theoretischen Gründen nachweisen.
\item[3)] Weit wichtiger als diese beiden Einwände ist der Umstand, dass P. das \textsc{Kritter}sche Gutachten falsch ausgelegt hat. Denn das, was P. unter dem Maximum der Witwenzahl verstand, und das, was man mit diesem Wort bezeichnet, wenn ein bestimmtes Normalverhältnis zwischen bestehenden Ehen und Witwen aufgestellt wird, und was auch in \textsc{Kritter}s Gutachten eigentlich gemeint ist,
\begin{center}\textit{sind zwei sehr verschiedene Dinge.}\end{center}
Hinzu kommt noch der ebenso wichtige Umstand,
\item[4)] dass die numerischen Ergebnisse, die für eine Gesellschaft von einem bestimmten, stets annähernd gleichbleibenden Umfang zulässig waren, wesentlich abgeändert werden müssen, wenn dieser Umfang sich bedeutend erweitert, wie dies oben schon ausführlich behandelt worden ist.
\end{list}
\par\medskip\centerline{\rule{.25\linewidth}{.3pt}}\medskip
